\section{自回归生成与模型输入}
\label{l12:sec:tokens}

\subsection{从条件分布到逐步生成}

大语言模型（Large Language Model，LLM）把文本生成表示为一个序列概率问题。分词器（tokenizer）先把文本转换为词元（token）的整数编号序列。词元可以对应一个字、一个子词、一段标点或若干字节；词元边界由分词规则决定，通常与自然语言中的词边界并不完全一致。设词表大小为 $V_{\mathrm{voc}}$，一个长度为 $N$ 的已知序列写为 $t_0,t_1,\ldots,t_{N-1}$，其中 $t_i\in\{0,\ldots,V_{\mathrm{voc}}-1\}$。

自回归（autoregressive）模型根据已有前缀预测下一个词元。对于一个完整序列，其联合概率可以按概率的链式法则分解为
\begin{equation}
  p(t_0,\ldots,t_{T-1})
  =\prod_{i=0}^{T-1}p(t_i\mid t_0,\ldots,t_{i-1}).
  \label{l12:eq:autoregressive}
\end{equation}
当 $i=0$ 时，条件为空，或由约定的起始标记提供。模型中的可训练参数共同确定这些条件分布。推理时固定参数，把系统输入、用户输入以及已经生成的词元共同作为上下文（context），反复计算下一词元的分布。

设完成最终归一化后，位置 $i$ 的表示向量为 $h_i\in\mathbb{R}^D$，$D$ 为模型宽度。词表投影把它变为未归一化分数（logits）$\ell_i=h_iW_U+b_U$，其中 $W_U\in\mathbb{R}^{D\times V_{\mathrm{voc}}}$，$b_U\in\mathbb{R}^{V_{\mathrm{voc}}}$ 为偏置向量。再经过软最大归一化（softmax）得到
\begin{equation}
  p(t_{i+1}=v\mid t_0,\ldots,t_i)
  =\frac{\exp(\ell_{i,v})}
  {\sum_{u=0}^{V_{\mathrm{voc}}-1}\exp(\ell_{i,u})}.
  \label{l12:eq:vocab-softmax}
\end{equation}
这里位置 $i$ 的输出对应位置 $i+1$ 的预测。这个错开一位的关系解释了因果注意力为什么允许位置 $i$ 读取它自身的输入。

生成过程还需要一个解码规则。贪心解码取分数最大的词元；随机采样则按照归一化后的分布选择词元。若采用温度 $\tau>0$，先把分数替换为 $\ell_i/\tau$，再计算分布。选出的词元追加到前缀后，进入下一次推理；遇到序列结束标记（end-of-sequence，EOS）、达到长度限制或满足其他既定停止条件时终止。词表分数的最大值位置在软最大归一化前后相同，因此只需贪心选择时可以直接对分数求最大值。

训练使用已知文本，可以同时构造许多位置的下一词元预测。推理时未来词元尚未确定，生成下一步的输入依赖当前步的选择。因此，同一已知前缀内部可以并行计算多个位置，连续的新词元生成仍存在步骤之间的依赖。

\subsection{分词、嵌入与位置信息}

分词器处理离散的符号边界与编号。字节对编码（Byte Pair Encoding，BPE）从较小的符号单元出发，根据语料中的共现统计反复合并高频相邻单元，形成词表和合并次序。推理时按照固定的规则把输入编码成编号。GPT-2 使用字节层面的 BPE；对同一段文字，具体编号必须与预训练模型使用的词表一致，不能通过临时给词语编号来替代。

词元嵌入（token embedding）把编号映射到连续向量。设嵌入矩阵为 $E\in\mathbb{R}^{V_{\mathrm{voc}}\times D}$，编号 $t_i$ 的嵌入就是第 $t_i$ 行 $E_{t_i,:}$。这等价于用独热向量左乘 $E$，实现时直接查表即可。编号本身只标识表中的行，编号之差不表示语义距离。嵌入向量之间的方向或内积关系可以携带训练得到的相似性，但一个坐标不必对应一个可独立命名的语义属性。

词表和嵌入表承担不同的任务。分词规则及合并次序通常在分词器构建阶段确定，随后保持固定；浮点嵌入参数通过语言模型的训练目标学习。推理既不重新学习词表，也不更新嵌入权重。GPT-2 小型模型的嵌入宽度为 $D=768$，所以一个整数编号被查表为一个含 $768$ 个浮点分量的向量。

词元嵌入本身不区分同一词元出现在哪个位置。位置信息需要另外提供。采用可训练的绝对位置嵌入时，设位置表为 $E^{\mathrm{pos}}\in\mathbb{R}^{N_{\max}\times D}$，其中 $N_{\max}$ 为模型支持的最大位置数。模型的初始输入为
\begin{equation}
  X^{(0)}_{i,:}=E_{t_i,:}+E^{\mathrm{pos}}_{i,:},
  \qquad 0\leq i<N\leq N_{\max}.
  \label{l12:eq:embedding}
\end{equation}
GPT-2 采用这种相加方式。嵌入矩阵 $X^{(0)}\in\mathbb{R}^{N\times D}$ 的行保留序列顺序，列表示特征。位置编码还可以采用其他形式；对这里的计算分析，关键条件是每个词元的位置在前缀扩展后保持一致。

查表与相加对每个位置读取并写入 $D$ 个分量，工作量为 $\Theta(ND)$。后续线性层常需 $\Theta(ND^2)$ 的算术工作，因此嵌入通常不是主体计算。不过，词表投影需要把宽度 $D$ 的表示映射到整个词表，涉及矩阵乘法，其开销应与输入查表区分。

\subsection{解码器结构与张量约定}

变换器（Transformer）架构利用注意力在不同位置之间传递信息。编码器把完整输入序列转换为上下文相关表示；在编码器与解码器组合的结构中，解码器还可以通过交叉注意力（cross-attention）读取编码器输出，并根据已生成前缀产生目标序列。仅含解码器的变换器（decoder-only Transformer）把输入提示和生成结果置于同一个序列中，通过因果自注意力（causal self-attention）建立依赖。GPT 系列的生成计算属于后一种结构。

设网络包含 $L$ 个变换器层。每层接收 $N\times D$ 的表示，输出形状保持为 $N\times D$，从而能够逐层堆叠。层内的多头注意力负责不同位置之间的信息混合，逐位置前馈网络负责每个位置内部的特征变换；残差连接和层归一化使这些变换能够稳定地组合。最终归一化和词表投影把最后一层表示转成下一词元分布。

\begin{table}[H]
\centering
\caption{张量与计数符号。矩阵公式默认省略批量维，向量按行参与乘法。}
\label{l12:tab:notation}
\small
\begin{tabularx}{\linewidth}{@{}L{3.4cm}Y L{3.5cm}@{}}
\toprule
符号 & 含义 & 形状或关系\\
\midrule
$B,N,L$ & 批量大小、上下文长度、层数 & 均为正整数\\
$D,h,d$ & 模型宽度、注意力头数、单头维度 & $D=hd$\\
$X$ & 某一层或子层的输入 & $N\times D$\\
$W_Q^{(a)},W_K^{(a)},W_V^{(a)}$ & 第 $a$ 个头的投影矩阵 & 每个为 $D\times d$\\
$Q^{(a)},K^{(a)},V^{(a)}$ & 单头的查询、键和值 & 每个为 $N\times d$\\
$S^{(a)},P^{(a)}$ & 单头分数矩阵、注意力权重 & 每个为 $N\times N$\\
$O^{(a)}$ & 单头输出 & $N\times d$\\
$D_{\mathrm{ff}}$ & 前馈网络的中间宽度 & 常取为 $D$ 的若干倍\\
$V_{\mathrm{voc}}$ & 词表大小 & 与值矩阵 $V^{(a)}$ 区分\\
\bottomrule
\end{tabularx}
\end{table}

例如，$N=4096$、$D=4096$、$h=32$ 时，$d=128$。输入仍然具有 $4096$ 个特征；每个注意力头把这 $4096$ 个特征投影到 $128$ 维空间。所有头的输出拼接后恢复到 $32\times128=4096$ 维。上下文长度和模型宽度恰好相等只是这个例子的取值，二者在定义和计算作用上独立。
